% !TEX root = ../../../tree_notes.tex
% =====================================================================
%  Chapter 13 review -- the 15-section revision contract (chapter_end.txt).
%  Revision-only: tables, one-liners, pointers back via \thmref/\probref/
%  \secref/\coderef/\defref/\exref/\figref/\propref.  No derivations.
%  Summit: CF 1076E Vasya and a Tree -- the per-depth lazy difference on
%  a rooted DFS, cur += V[depth] on entry, deferred V[depth+d+1] -= x per
%  query, undo on exit.  Sample verified live (Luogu mirror P3674 / CF
%  statement): n=5, edges 1-2 1-3 2-4 2-5, queries (1,1,1) (2,0,10)
%  (4,10,100) -> "1 11 1 100 0", hand-traced below.
%  Exports prob:c10-summit, lst:c10-summit, lst:c10-kit (this chapter's
%  three remaining forward-usable labels).
% =====================================================================

\chaptersection{Chapter 13 Summary}

One sentence: \emph{a fixed root turns every tree into a fold -- one tuple
per node, combined over children by a rule that names the DP; the four
shapes this chapter drills are scalar fold, in/out pair, arm-at-peak
harvest, and depth-indexed histogram, and every problem below is one of
them with the right $\Phi$.} Chapter~12 measured a tree; Chapter~13 makes
it compute. The DFS skeleton \coderef{c09-dfs3} returns verbatim; what
changes is the state you carry on each frame and the combiner that folds
children into the parent. Every entry below is a pointer, not a
re-derivation.

\begin{center}
\renewcommand{\arraystretch}{1.25}
\begin{tabular}{@{}p{0.24\linewidth} p{0.48\linewidth} p{0.20\linewidth}@{}}
\toprule
\textbf{Result} & \textbf{Statement (with conditions)} & \textbf{Where} \\
\midrule
Subtree fold & On a rooted tree, a state $\mathcal S(v)=\Phi(\text{self}(v),\{\mathcal S(c):c\in\children(v)\})$ computed in one DFS is a \emph{template}; LC~333's $(\mathit{isBST},\mathit{lo},\mathit{hi},\mathit{sz},\mathit{best})$ tuple is one instantiation. & \defref{c10-subtree-state}, \thmref{c10-subtree-template} \\
Constant state & If $|\mathcal S(v)|$ is a $k$-coordinate tuple independent of $\deg(v)$, the fold costs $\bigO{n}$. & \thmref{c10-subtree-template} \\
Two-state in/out & $\mathit{dp}_{\mathit{in}}[v]=w(v)+\sum_c\mathit{dp}_{\mathit{out}}[c]$, $\mathit{dp}_{\mathit{out}}[v]=\sum_c\max(\mathit{dp}_{\mathit{in}}[c],\mathit{dp}_{\mathit{out}}[c])$ solves weighted MIS on a tree; vertex cover is the complement by K{\H{o}}nig. & \defref{c10-inout}, \thmref{c10-mis-tree} \\
Harvest at peak & Every path has a unique topmost node $\ell=\mathrm{peak}(u,w)=\mathrm{LCA}(u,w)$; any additive path functional factors as $\mathrm{combine}(\mathrm{arm}(u,\ell),\mathrm{arm}(w,\ell))$ and sums over $\ell$. & \propref{c10-pairtop}, \thmref{c10-pairtop-count} \\
Distance-$k$ count & $\sum_{\{u,w\}}[d(u,w)=k]=\sum_\ell\sum_{\text{arms at }\ell}[i+j=k]$; the seed $\mathit{H}[u][0]=1$ is the length-$0$ arm (one endpoint is $\ell$). & \thmref{c10-pairtop-count}, \exref{c10-k2} \\
Depth histogram & State $\mathit{H}_v[j]$ at $v$ folds $\Phi$ over $v$'s subtree at relative depth $j$; child merge shifts $+1$ and costs $\bigO{k}$ per edge. & \defref{c10-dist-hist} \\
Invertible accum. & Offline per-node, per-depth queries with an invertible $\Phi$ cost $\bigO{n+m}$ via a snapshot-then-read pair at each $v$'s entry/exit -- \emph{no} restore on exit. & \thmref{c10-depth-xor} \\
\bottomrule
\end{tabular}
\end{center}

\medskip
{\small\textbf{Algorithms from this chapter -- one DFS, one combiner, one answer per node.}}
\begin{center}
\small
\renewcommand{\arraystretch}{1.2}
\setlength{\tabcolsep}{3pt}
\begin{tabular}{@{}p{0.19\linewidth} l l p{0.29\linewidth} p{0.155\linewidth}@{}}
\toprule
\textbf{Algorithm} & \textbf{Time} & \textbf{Memory} & \textbf{When to use (deciding phrase)} & \textbf{Code} \\
\midrule
Subtree tuple fold
  & $\bigO{n}$ & $\bigO{n}$
  & ``an answer per node computed bottom-up from a fixed-size tuple''
  & \coderef{c10-fold} \\
In/out pair (MIS)
  & $\bigO{n}$ & $\bigO{n}$
  & ``take or leave $v$, forced on children''
  & \coderef{c10-inout} \\
Peak-harvest histogram
  & $\bigO{nk}$ & $\bigO{nk}$
  & ``count/combine over all pairs at distance $k$''
  & \coderef{c10-peak} \\
Distance histogram (Shape A)
  & $\bigO{nk}$ & $\bigO{nk}$
  & ``bottom-up per-node depth table''
  & \coderef{c10-distk-hist} \\
Backtrack DFS + prefix XOR (Shape B)
  & $\bigO{n+m}$ & $\bigO{n+m}$
  & ``offline per-node, per-depth read with an invertible $\Phi$''
  & \coderef{c10-depth-hist} \\
Per-depth lazy difference (Summit)
  & $\bigO{n+m}$ & $\bigO{n}$
  & ``$m$ range-by-depth subtree updates, final value at every node''
  & \coderef{c10-summit} \\
\bottomrule
\end{tabular}
\end{center}

\chaptersection{Master Decision Flowchart}

Post-study ladder: walk top to bottom, stop at the first branch that
matches. It decides \emph{which of this chapter's four DP shapes} the
problem is asking for -- and hands off what this chapter does \emph{not}
own: the answer for every root (Chapter~15's rerooting) and per-pair
\emph{query} distance (Chapter~16's LCA).

\begin{lstlisting}[language={}, basicstyle=\ttfamily\small, numbers=none, frame=single]
  You are handed a rooted tree and asked for a value.  Where does it live?
  |
  |-- one scalar per node, folded from children by a fixed rule?
  |     +-- tuple state, one DFS: [10.1, LC 333]
  |
  |-- a binary decision at every node (take / leave, colour / not)?
  |     +-- in/out pair, sum-max recurrence:     [10.2, CF 1529C]
  |        (weighted MIS on trees, vertex cover dual)
  |
  |-- a value defined on PATHS (not nodes), summed over all pairs?
  |     +-- harvest at the peak of each path:    [10.3, CF 161D]
  |        sum_{u<w} F(u,w) = sum_l sum_{arms at l} combine(.,.)
  |        (Ch.12's thm:c09-dpjoin with a general Phi)
  |
  |-- a value indexed by a distance k in [0, n]?
  |     +-- per-node depth histogram, merge children: [10.4 Shape A]
  |        cost O(nk); Ch.14 (c27) analyses the budget.
  |
  |-- OFFLINE per-node, per-depth aggregate with an invertible combiner?
  |     +-- backtracking DFS with a monotone prefix accumulator:
  |        [10.4 Shape B, CF 570D]  cost O(n+m), never restore X[d].
  |
  |-- SUBTREE update queries with a per-depth budget (v, d, x)?
  |     +-- per-depth lazy difference V[d]; enter: cur += V[dep],
  |        apply queries, recurse, undo.        [Summit CF 1076E]
  |
  |-- "for every root"?
  |     +-- NOT this chapter -- Ch.15's rerooting theory (c20).
  |
  |-- d(u, v) for MANY arbitrary pairs, ONLINE?
  |     +-- NOT this chapter -- Ch.16's LCA arithmetic (c11).
  |
  +-- knapsack budget with a per-node capacity w?
        +-- NOT this chapter -- Ch.14's formal analysis (c27);
            here we preview the sum-of-subtree-size bound only.
\end{lstlisting}

\chaptersection{Common Mistakes Checklist}

\begin{itemize}[nosep, leftmargin=1.4em]
  \item \textbf{State too small.} A tuple that drops the coordinate that
        disambiguates two children loses information the fold cannot
        recover. LC~333's $(\mathit{isBST},\mathit{lo},\mathit{hi},
        \mathit{sz},\mathit{best})$ is the smallest honest budget: cutting
        $\mathit{lo}$/$\mathit{hi}$ makes the BST check myopic at the
        grandparent boundary (\probref{c10-subtree-solved},
        \defref{c10-subtree-state}).
  \item \textbf{Forgetting the parent-side piece.} A subtree DP answers
        for the subtree \emph{rooted at}~$v$; if the problem asks about
        ``paths through $v$'' you must include the arm going up through
        $v$'s parent. Ch.13 does not solve the every-root version --
        that is Chapter~15's rerooting (\chref{c20}); recognising the
        need is the boundary check.
  \item \textbf{Arm-at-peak double count.} In \coderef{c10-peak},
        \emph{harvest before you merge}. If you merge then harvest,
        same-child pairs get counted at their parent's peak -- the WA
        signature is ``answer too big by a factor of two at every
        branching node'' (\thmref{c10-pairtop-count},
        \probref{c10-harvest-solved}).
  \item \textbf{Dropping the seed $\mathit{H}[u][0]=1$.} The row where
        $\ell$ \emph{is} one of the endpoints contributes $\sum_c
        \mathit{H}_c[k-1]$ at every distance-$k$ query; forgetting it is
        the top WA cause on CF~161D (\exref{c10-harvest-run},
        \exref{c10-k2}).
  \item \textbf{Distance-$k$ off-by-one.} $\mathit{H}_v[j]$ measures the
        distance from $v$; a child's $\mathit{H}_c[j]$ lands at
        $\mathit{H}_v[j+1]$ after the shift. Writing
        $\mathit{H}_v[j] += \mathit{H}_c[j]$ (no $+1$) silently drops
        every pair that crosses a real edge (\coderef{c10-peak},
        \coderef{c10-distk-hist}).
  \item \textbf{Restoring the accumulator in Shape~B.}
        \thmref{c10-depth-xor}'s power is that $\mathit{X}[d]$ is
        \emph{never} restored: the difference
        $\mathit{X}[h]_{\mathrm{out}}\ominus\mathit{X}[h]_{\mathrm{in}}$
        is the subtree aggregate. XOR-restoring on $u$'s exit collapses
        the difference to $0$ for every internal query
        (\probref{c10-distk-solved}, 13.4's Pitfalls).
  \item \textbf{\lstinline{int} overflow on CF~1529C and CF~1076E.}
        Weighted MIS sums reach $10^{14}$ at $n=10^{5}$, $a\le 10^{9}$;
        the Summit's $x$-sums reach $n \cdot \max x = 3\cdot10^{14}$.
        Every accumulator is \lstinline{long long} -- this is Ch.12's
        trap re-appearing inside a fold rather than a shift
        (\probref{c10-inout-solved}, \probref{c10-summit}).
  \item \textbf{Recursion depth on a legal chain.} $n=5\cdot10^{5}$
        blows a naive recursive DFS; Ch.13's solved listings use the
        \emph{iterative two-phase stack} introduced in Ch.12's
        \coderef{c09-iter}, applied to the DP shape (\secref{c10-subtrees},
        \probref{c10-distk-solved}).
\end{itemize}

\chaptersection{WA Clinic}

Verdicts are symptoms; the causes are the short list this chapter
drilled. Diagnose top-down -- every leaf names the training problem.

\begin{figure}[H]
\centering
\begin{tikzpicture}[diagram,
    v/.style={boxnode, fill=primaryblue!12, draw=primaryblue, text width=22mm, align=center},
    b/.style={boxnode, fill=difforange!12, draw=difforange, text width=12mm},
    l/.style={boxnode, fill=successgreen!10, draw=successgreen, text width=58mm,
              font=\scriptsize, align=left}]
  \node[v] (verdict) {\textbf{Verdict on the judge}};
  \node[b, below left=8mm and 6mm of verdict] (wa) {\textbf{WA}};
  \node[b, below right=1mm of verdict] (tle) {\textbf{TLE}};
  \node[b, below left=8mm and 1mm of verdict, xshift=30mm] (re) {\textbf{RE}};
  \node[b, below right=8mm and 6mm of verdict] (mle) {\textbf{MLE}};
  \node[l, below right=1mm and 2mm of wa.south] (wa1)
    {merge-then-harvest at the peak double-counts same-child pairs (\probref{c10-harvest-solved});
     dropping the seed $\mathit{H}[u][0]=1$ skips the length-$0$-arm rows (\exref{c10-k2});
     in-state using $\max(\mathit{dp}_{\mathit{in}}[c],\mathit{dp}_{\mathit{out}}[c])$ violates independence (\thmref{c10-mis-tree});
     restoring $\mathit{X}[d]$ in Shape B zeros the difference (\probref{c10-distk-solved})};
  \node[l, below right=1mm and 2mm of tle.south] (tle1)
    {$\bigO{nk}$ histogram merge with $k$ up to $n$ -- use Shape B if $\Phi$ is invertible (\thmref{c10-depth-xor});
     brute-force pair enumeration $\bigO{n^2}$ where peak-harvest is $\bigO{nk}$ (\propref{c10-pairtop})};
  \node[l, below right=1mm and 2mm of re.south] (re1)
    {recursive DFS on a legal $5\cdot10^{5}$-node chain -- the stack, not the answer, dies;
     use \coderef{c09-iter}'s two-phase form inside \coderef{c10-fold} (\secref{c10-subtrees})};
  \node[l, below right=1mm and 2mm of mle.south] (mle1)
    {$\bigO{nk}$ histogram per node with $k$ unbounded -- small-to-large (Ch.17's \chref{c28}) is the fix;
     $\mathit{H}$ declared as a dense matrix sized $n\times n$ instead of $n\times(k+1)$};
  \draw[flow] (verdict) -- (wa); \draw[flow] (verdict) -- (tle);
  \draw[flow] (verdict) -- (re); \draw[flow] (verdict) -- (mle);
  \draw[flow] (wa) -- (wa1); \draw[flow] (tle) -- (tle1);
  \draw[flow] (re) -- (re1); \draw[flow] (mle) -- (mle1);
\end{tikzpicture}
\caption*{Chapter~13 verdict tree -- only the failures this chapter's
problems can produce.}
\end{figure}

\medskip
\noindent Backup table (same content, terse):
\begin{center}
\small
\renewcommand{\arraystretch}{1.2}
\begin{tabular}{@{}llp{0.39\linewidth}l@{}}
\toprule
\textbf{Verdict} & \textbf{Usual cause} & \textbf{One-line fix} & \textbf{Trained in} \\
\midrule
WA  & merge-then-harvest       & \code{harvest then shift-merge}       & \probref{c10-harvest-solved} \\
WA  & no seed $\mathit{H}[u][0]=1$ & write the seed before the child loop & \exref{c10-k2} \\
WA  & \code{dp\_in} uses $\max$    & replace with \code{dp\_out[c]} (unconditional) & \thmref{c10-mis-tree} \\
WA  & Shape B restores $\mathit{X}$ & remove the restore line, keep snapshot-diff & \probref{c10-distk-solved} \\
TLE & $\bigO{n^2}$ pair loop        & $\bigO{nk}$ histogram merge             & \coderef{c10-peak} \\
RE  & recursion on a chain           & \coderef{c09-iter} two-phase stack       & \secref{c10-subtrees} \\
\bottomrule
\end{tabular}
\end{center}

\chaptersection{Quick Reference Card}

\begin{quickref}[Chapter 13 at a glance]
\textbf{The DFS skeleton never changes.} \coderef{c09-dfs3} returns; only
the tuple $\mathcal S(v)$ and the combiner $\Phi$ move
(\defref{c10-subtree-state}, \thmref{c10-subtree-template}). On a deep
chain, the swap to \coderef{c09-iter} is mechanical.

\medskip
\textbf{Four shapes, one decision.} Scalar fold (13.1); in/out pair for a
binary decision per node (13.2); arm-at-peak when the objective is
\emph{on paths} (13.3); per-depth histogram when the objective is
\emph{indexed by distance} (13.4). Naming the shape is 80\% of the solve
(the taxonomy chart in the chapter opener).

\medskip
\textbf{Harvest before merge, seed with 1.} In \coderef{c10-peak}, the
pair $\{u,w\}$ is counted \emph{once}, at $\mathrm{peak}(u,w)$. Merging
first double-counts same-child pairs; missing the seed
$\mathit{H}[u][0]=1$ drops the length-$0$-arm rows
(\propref{c10-pairtop}, \thmref{c10-pairtop-count}, \exref{c10-k2}).

\medskip
\textbf{Invert $\Phi$ for free $\bigO{n+m}$.} If the combiner is XOR (or
any group operation), \thmref{c10-depth-xor} beats the $\bigO{nk}$
histogram merge by an order of magnitude -- snapshot at entry, read the
difference at exit, \emph{never restore}. This is CF~570D's whole trick
and the shape the Summit generalises to a lazy difference array.

\medskip
\textbf{Know your boundaries.} Answer for every root is Ch.15
(\chref{c20}); per-pair query distance is Ch.16 (\chref{c11}); knapsack
budget analysis is Ch.14 (\chref{c27}); small-to-large is Ch.17
(\chref{c28}); centroid decomposition is Ch.24 (\chref{c18}). This
chapter owns the \emph{fixed-root DP cookbook}.
\end{quickref}

\begin{lstlisting}[caption={Chapter 13 code kit: the tuple fold, the in/out pair, the peak-harvest merge, the two depth-histogram shapes, and the Summit's per-depth lazy difference}, label={lst:c10-kit}]
#include <bits/stdc++.h>
using namespace std;
int n;  vector<vector<int>> adj;

// (1) Subtree tuple fold:  S(v) = Phi(self(v), { S(c) }).           [10.1]  lst:c10-fold
struct S { bool ok; int lo, hi, sz, best; };
S fold(int u, int p, const vector<int>& val) {
    S r{true, val[u], val[u], 1, 1};
    for (int v : adj[u]) if (v != p) {
        S c = fold(v, u, val);                    // parent guard = the rooting
        r.lo = min(r.lo, c.lo);  r.hi = max(r.hi, c.hi);
        r.sz += c.sz;  r.best = max(r.best, c.best);
        r.ok = r.ok && c.ok;                      // fold is coordinate-wise
    }
    // ... (LC 333 BST gate updates best to r.sz when r.ok)
    return r;
}

// (2) In/out pair (MIS): take v forces children out; skip v leaves them free. [10.2] lst:c10-inout
vector<long long> dp_in, dp_out;
void mis(int u, int p, const vector<long long>& w) {
    dp_in[u] = w[u];  dp_out[u] = 0;
    for (int v : adj[u]) if (v != p) {
        mis(v, u, w);
        dp_in[u]  += dp_out[v];                    // unconditional out
        dp_out[u] += max(dp_in[v], dp_out[v]);     // free choice
    }
}

// (3) Peak-harvest histogram, distance k: harvest BEFORE merge, seed H[u][0]=1. [10.3] lst:c10-peak
int K;  vector<vector<int>> H;  long long ans;
void peak(int u, int p) {
    H[u].assign(K + 1, 0);
    H[u][0] = 1;                                   // length-0 arm (u itself)
    for (int v : adj[u]) if (v != p) {
        peak(v, u);
        for (int i = 0; i <= K - 1; ++i)           // HARVEST first
            ans += (long long)H[u][i] * H[v][K - 1 - i];
        for (int j = K; j >= 1; --j)               // SHIFT-MERGE after
            H[u][j] += H[v][j - 1];
    }
}

// (4) Shape B: backtracking DFS with a per-depth XOR accumulator, never restored. [10.4] lst:c10-depth-hist
vector<int> X;                                     // 26-bit masks, indexed by depth
void xor_hist(int u, int p, int d, const vector<string>& args) {
    // (a) snapshot X[h] per query bucketed at u, before adding u's own bit
    // (b) X[d] ^= (1 << (s[u]-'a'));
    // (c) recurse children at depth d+1;
    // (d) at u's exit, answer each query as (X[h] ^ in_val[id]) popcount <= 1
    // (e) NO RESTORE -- the difference out ^ in is the subtree aggregate.
}

// (5) Summit per-depth lazy difference: cur += V[dep], undo on exit. [Summit 1076E] lst:c10-summit
long long cur;  vector<long long> V;               // V[d] = deferred delta for entering depth d
void vasya(int u, int p, int d, const vector<vector<pair<int,int>>>& qs, vector<long long>& ans) {
    long long cur0 = cur, V0 = V[d];
    cur += V[d];  V[d] = 0;                        // consume the deferred delta
    for (auto& [dq, x] : qs[u]) { cur += x; V[d + dq + 1] -= x; }  // stop applying past d+dq
    ans[u] = cur;
    for (int v : adj[u]) if (v != p) vasya(v, u, d + 1, qs, ans);
    cur = cur0;  V[d] = V0;                        // undo on exit
}
\end{lstlisting}

\chaptersection{Write-From-Memory Drills}

Five signature skeletons, each lifted from a listing already in the
chapter. Cover the marked line, reproduce it from memory, then check the
answer one line below.

\begin{lstlisting}[caption={Drill 1 -- the subtree tuple fold (from \coderef{c10-fold})}]
S fold(int u, int p) {
    S r = /* FILL: self(v) */;                     // leaf initialisation
    for (int v : adj[u]) if (v != p) {             // parent guard = the rooting
        S c = fold(v, u);
        r = /* FILL: Phi(r, c) */;                 // coordinate-wise combine
    }
    return r;                                       // answer lives in r.best
}
\end{lstlisting}
\noindent{\scriptsize Key: the fold is a monoid over children -- order
does not matter, so the loop can go in any child sequence. LC~333's
$(\mathit{ok},\mathit{lo},\mathit{hi},\mathit{sz},\mathit{best})$ is a
five-coordinate instance; a smaller tuple loses information at the
grandparent boundary (\defref{c10-subtree-state}, \probref{c10-subtree-solved}).}

\begin{lstlisting}[caption={Drill 2 -- the in/out pair recurrence (from \coderef{c10-inout})}]
void mis(int u, int p) {
    dp_in[u] = w[u];  dp_out[u] = 0;
    for (int v : adj[u]) if (v != p) {
        mis(v, u);
        dp_in[u]  += /* FILL: dp_out[v] */;                 // unconditional
        dp_out[u] += /* FILL: max(dp_in[v], dp_out[v]) */;   // free choice
    }
}
\end{lstlisting}
\noindent{\scriptsize Key: the in-state \emph{forces} every child out --
using $\max$ there is Trap~1 of \secref{c10-twostate}. The out-state is
unconstrained and takes $\max$ per child, \emph{summed} over children
(the outer $\sum$ is what makes the fold additive, not competitive).
Weighted MIS on a tree, K{\H{o}}nig's vertex-cover dual, and CF~1529C's
convexity-at-endpoints all read off this one pair (\thmref{c10-mis-tree}).}

\begin{lstlisting}[caption={Drill 3 -- harvest at the peak (from \coderef{c10-peak})}]
void peak(int u, int p) {
    H[u].assign(K + 1, 0);
    /* FILL: H[u][0] = 1 */;                       // seed the length-0 arm
    for (int v : adj[u]) if (v != p) {
        peak(v, u);
        for (int i = 0; i <= K - 1; ++i)
            ans += (long long)/* FILL: H[u][i] * H[v][K-1-i] */;   // HARVEST
        for (int j = K; j >= 1; --j)
            H[u][j] += H[v][j - 1];                // SHIFT-MERGE (+1 depth)
    }
}
\end{lstlisting}
\noindent{\scriptsize Key: \emph{harvest before merge}. Once $v$'s
histogram is shifted into $u$'s, same-child pairs are indistinguishable
from cross-child ones and get counted at $u$'s parent's peak -- the top
WA cause on CF~161D. The seed $\mathit{H}[u][0]=1$ is the arm where one
endpoint \emph{is} the peak; Ch.16's \exref{c10-k2} forward reference
relies on this row (\propref{c10-pairtop}, \thmref{c10-pairtop-count}).}

\begin{lstlisting}[caption={Drill 4 -- Shape B: per-depth XOR accumulator (from \coderef{c10-depth-hist})}]
void dfs(int u, int p, int d) {
    for (auto& [h, id] : q[u]) in_val[id] = /* FILL: X[h] */;    // snapshot
    X[d] ^= (1 << (s[u] - 'a'));                                  // add u's own bit
    for (int v : adj[u]) if (v != p) dfs(v, u, d + 1);
    for (auto& [h, id] : q[u])
        ans[id] = (popcount(X[h] ^ in_val[id]) <= 1);             // subtree aggregate
    /* FILL: NO restore of X[d] */;
}
\end{lstlisting}
\noindent{\scriptsize Key: the accumulator $\mathit{X}[\cdot]$ is
\emph{monotone}. Restoring on exit collapses the difference to zero
because children also restored their bits -- \thmref{c10-depth-xor}
requires a group operation, and $\oplus$ is its own inverse. Snapshot
must be \emph{before} $u$'s own bit, not after, or you exclude $u$ from
its own subtree query (\probref{c10-distk-solved}).}

\begin{lstlisting}[caption={Drill 5 -- Summit: per-depth lazy difference (from \coderef{c10-summit})}]
void vasya(int u, int p, int d) {
    long long cur0 = cur, V0 = V[d];
    cur += V[d];  V[d] = 0;                        // consume the deferred delta
    for (auto& [dq, x] : qs[u]) { cur += x; V[d + dq + 1] -= x; }
    ans[u] = cur;
    for (int v : adj[v]) if (v != p) vasya(v, u, d + 1);
    /* FILL: cur = cur0; V[d] = V0; */             // undo on exit -- this is the restore
}
\end{lstlisting}
\noindent{\scriptsize Key: \emph{here} the restore is right. The
accumulator is a per-node value $\mathit{cur}$, not a per-depth
monotone aggregate; the deferred $\mathit{V}[d{+}dq{+}1] \mathrel{-}= x$
is applied when a \emph{deeper} node enters, and siblings must not see
the delta. Compare with Drill 4: whether to restore depends on the
question Shape B is answering (\probref{c10-summit}).}

\chaptersection{Complete Problem Checklist}

Every problem this chapter offers -- 4 fully solved in-section, 4
practice, one Summit. Tick \emph{only} after an unaided solve; a hinted
attempt counts as half. The ten mixed problems are listed in their own
station. Cite by name and number -- URLs appear once, in the solved and
practice entries.

\begin{center}
\footnotesize
\renewcommand{\arraystretch}{1.25}
\begin{tabular}{@{}c p{0.19\linewidth} l p{0.20\linewidth} p{0.235\linewidth} c@{}}
\toprule
\textbf{\#} & \textbf{Problem} & \textbf{Source} & \textbf{Key technique} & \textbf{What to look for} & \textbf{Done?} \\
\midrule
\multicolumn{6}{l}{\textbf{Section 10.1 -- state per node, tuple fold}} \\
W13.1.1 & Largest BST Subtree & LC 333 & \lv{3} five-coordinate tuple & BST flag + lo/hi propagated up; sz used only to size a valid BST & $\square$ \\
P13.1.2 & Count Ways Where You Can Find the Best Robbers & LC 2581 & \lv{4} weighted MIS variant with parity + counting & Ch.2 owns the URL; re-solve cold with the in/out pair and a count array & $\square$ \\
\midrule
\multicolumn{6}{l}{\textbf{Section 10.2 -- in/out pair (MIS / vertex cover)}} \\
W13.2.1 & Parsa's Humongous Tree & CF 1529C & \lv{3} convexity-at-endpoints + in/out pair & Every $a_v$ is either $\ell_v$ or $r_v$; per-edge weight matrix; \lstinline{long long} & $\square$ \\
P13.2.2 & Marginal Cost & CF 1010D & \lv{4} scalar fold, one state per node & Ch.\,13.2's fold skeleton, single state; recognising state budget is the drill & $\square$ \\
\midrule
\multicolumn{6}{l}{\textbf{Section 10.3 -- harvest at the peak}} \\
W13.3.1 & Distance in Tree & CF 161D & \lv{4} histogram merge, seed $\mathit{H}[u][0]=1$ & Harvest \emph{before} merge; $\bigO{nk}$; the seed row is what Ch.16's degenerate census rides on & $\square$ \\
P13.3.2 & Child and Puzzle & CF 839B & \lv{3} top-two-arms predicate at peak & Three-coordinate tuple $(\mathit{top}_1,\mathit{top}_2,\mathit{ok})$ -- the Ch.12 diameter geometry with a budget & $\square$ \\
\midrule
\multicolumn{6}{l}{\textbf{Section 10.4 -- distance-indexed DP, two shapes}} \\
W13.4.1 & Tree Requests & CF 570D & \lv{4} backtracking DFS, monotone XOR accumulator & Never restore $\mathit{X}[\cdot]$; snapshot per query at entry, read difference at exit & $\square$ \\
P13.4.2 & Directing the Edges & CF 1385E & \lv{3} scalar fold with an in/out pair of flip counts & Root at $1$, one DFS for $\mathit{down}[v]$; $\min_v$ over every candidate root & $\square$ \\
\midrule
\multicolumn{6}{l}{\textbf{Summit (chapter end)}} \\
S13 & Vasya and a Tree & CF 1076E & \lv{4} per-depth lazy difference on a DFS & cur += V[dep], V[dep+dq+1] -= x, undo on exit; \lstinline{long long} throughout & $\square$ \\
\bottomrule
\end{tabular}
\end{center}

\chaptersection{Recommended Study Order}

Five focused days. Day~1 rebuilds the tuple fold and internalises the
``fixed-size state'' budget; day~2 doubles the state to a pair and
drills the convexity reduction; day~3 harvests paths at the peak; day~4
runs the two Shape-A/Shape-B variants of the depth-indexed state; day~5
is the Summit. Solve in this order -- each day's tool is the previous
day's output.

\begin{center}
\small
\renewcommand{\arraystretch}{1.25}
\begin{tabular}{@{}llp{0.55\linewidth}@{}}
\toprule
\textbf{Day} & \textbf{Focus} & \textbf{Tasks} \\
\midrule
1 & Read + do 10.1 & trace $T$'s tuple fold at $\mathit{val}=[3,1,4,1,5,9,2,6,5,3]$; write \coderef{c10-fold} cold; LC 333; LC 2581 as a preview of the pair fold \\
2 & Read + do 10.2 & verify \thmref{c10-mis-tree}'s table $\mathit{dp}_{\mathit{in}}/\mathit{dp}_{\mathit{out}}$ on $T$ (\exref{c10-mis-run}); CF 1529C; CF 1010D as a scalar budget check \\
3 & Read + do 10.3 & walk the k=2 census on $T$ by hand (\exref{c10-harvest-run}); CF 161D; CF 839B as a top-two-arms predicate \\
4 & Read + do 10.4 & do CF 570D twice -- once with Shape A (histogram merge) to see the cost, once with Shape B; CF 1385E as a scalar exit-time read \\
5 & Summit + revise & CF 1076E with the per-depth difference array (\probref{c10-summit}); walk the decision flowchart on five new problem statements \\
\bottomrule
\end{tabular}
\end{center}

\chaptersection{Difficulty Progression}

Tiers are competency gates -- each row names the \emph{skill} the tier
unlocks, not just a rating bucket.

\begin{center}
\small
\renewcommand{\arraystretch}{1.3}
\begin{tabular}{@{}lp{0.36\linewidth}p{0.34\linewidth}@{}}
\toprule
\textbf{Tier} & \textbf{Problems} & \textbf{What passing this tier proves} \\
\midrule
$\lv{3}$ \ $\bigstar\bigstar\bigstar$
  & LC 333, CF 1529C, CF 839B, CF 1385E
  & you fix a root, name a constant-size state, and fold it through a child loop without re-deriving Ch.~1's recursion \\
$\lv{4}$ \ $\bigstar\bigstar\bigstar\bigstar$
  & LC 2581, CF 161D, CF 570D, CF 1010D, CF 1076E (Summit)
  & you index the state (pair, arm, histogram, depth) and match the combiner to the objective; you know when to restore and when not to \\
$\lv{5}$ \ $\bigstar\bigstar\bigstar\bigstar\bigstar$
  & CSES 2103/2104 (LCA previews), CF 1324F / 763B (rerooting previews)
  & you \emph{recognise} the shape and hand off to Ch.15 / Ch.16 when it goes beyond one fixed-root DP \\
\bottomrule
\end{tabular}
\end{center}

\chaptersection{Reduction Gallery}

Every DP is ``fix a root, name a state, fold over children.'' Each line
collapses a family onto a listing already in the chapter -- no new
theory.

\begin{itemize}[nosep, leftmargin=1.4em]
  \item ``largest subtree that is a BST'' (LC~333) $\longrightarrow$
        five-coordinate tuple fold; $\Phi$ reads
        ``min, max, size, ok, best'' over each child's tuple
        (\probref{c10-subtree-solved}, \coderef{c10-fold}).
  \item ``max total value of non-adjacent nodes'' (LC 337 on binary
        trees, CF~1529C on general trees) $\longrightarrow$ in/out pair
        (\thmref{c10-mis-tree}, \coderef{c10-inout}); Ch.2's
        $\mathit{rob}/\mathit{skip}$ is the $\deg\le 2$ special case.
  \item ``count paths of length $k$'' (CF 161D, CSES 2081)
        $\longrightarrow$ peak-harvest with a depth-indexed histogram
        (\propref{c10-pairtop}, \coderef{c10-peak}); equality becomes a
        range test, everything else is unchanged.
  \item ``is the character multiset at depth $h$ in $v$'s subtree a
        palindrome?'' (CF 570D) $\longrightarrow$ backtracking DFS with
        a monotone XOR accumulator (\thmref{c10-depth-xor},
        \coderef{c10-depth-hist}).
  \item ``direct every edge away from some root'' (CF 1385E)
        $\longrightarrow$ scalar pair $(\mathit{down}[v],\mathit{up}[v])$
        on the same fold (\secref{c10-subtrees}, \coderef{c10-fold});
        the $\min_v$ over roots is not a reroot -- one scalar table +
        one forward walk.
  \item ``$m$ subtree-depth-range updates, final value per node''
        (CF 1076E) $\longrightarrow$ per-depth lazy difference on a DFS
        (\probref{c10-summit}, \coderef{c10-summit}); Shape B's
        snapshot idea, promoted to $\Phi=+$.
\end{itemize}

\chaptersection{Mixed Problem Set}

These refuse to say ``DP''; decide \emph{which quantity lives at each
node and what $\Phi$ is} before opening the code. A hint tells you the
shape of the answer, never the tool.

\begin{mixedset}
  \practice{Subordinates}{CSES 1126}
    {https://cses.fi/problemset/task/1126}{CSES · L3}
    {Given the direct boss of every employee (a rooted tree of $n$ nodes), count subordinates in each node's subtree. One bottom-up fold: $\mathit{sz}[v] = 1 + \sum_c \mathit{sz}[c]$ -- \secref{c10-subtrees}'s simplest instantiation, and Ch.12's $\size$ re-rolled as a ``who-below-me'' table.}
  \practice{Company Queries II}{CSES 2103}
    {https://cses.fi/problemset/task/2103}{CSES · L4}
    {Offline $k$-th ancestor per node -- this chapter's parent-guarded DFS builds the parent table; the binary-lifting overlay is Ch.16's (\chref{c11}). Do the DFS half here; the $\bigO{\log n}$ lift is a forward pointer, not a Ch.13 solve.}
  \practice{Company Queries III}{CSES 2104}
    {https://cses.fi/problemset/task/2104}{CSES · L4}
    {LCA + subtree queries on a rooted tree; the DP here is \emph{only} the depth/parent precompute. Same forward-pointer discipline as 2103 -- Ch.13 hands you the fold, Ch.16 handles the query structure (\chref{c11}).}
  \practice{Tree (Two-Colour Chain)}{Codeforces 55C}
    {https://codeforces.com/problemset/problem/55/C}{CF · L4}
    {A path-like tree needs to be coloured so that any three consecutive nodes have at least two colours. The DP state is $(\mathit{colour}(v),\mathit{colour}(\mathit{parent}(v)))$, a two-coordinate pair -- \thmref{c10-mis-tree}'s shape with a different $\Phi$ (count instead of $\max$). Root at~$1$.}
  \practice{GCD Counting}{Codeforces 1101D}
    {https://codeforces.com/problemset/problem/1101/D}{CF 1700 · L4}
    {Longest path where every node value is divisible by~$g$. The peak-harvest geometry of \propref{c10-pairtop} with $F(u,w)=[g\mid \gcd(\text{nodes on path})] \cdot \mathit{len}$; the arm is the longest downward chain divisible by~$g$, and the combine is ``top two arms at the peak plus one'' (\coderef{c10-peak}).}
  \practice{Maximum White Subtree}{Codeforces 1324F}
    {https://codeforces.com/problemset/problem/1324/F}{CF 1800 · L4}
    {Colour each node black/white; for every root $r$, compute (black $-$ white) in $r$'s rooted tree, then promote to ``best connected monochrome piece.'' The single-root version is \secref{c10-twostate}'s fold; the every-root version is Ch.15's \chref{c20} reroot territory. Do it here with $r=1$ fixed and \emph{flag} the extra work you'd need for all roots -- that flag is your Ch.15 entrance ticket.}
  \practice{Popular Vote}{Codeforces 763B}
    {https://codeforces.com/problemset/problem/763/B}{CF 2000 · L5}
    {Remove a connected subtree so that the remaining forest has a strict majority value; find the smallest removable subtree. The predicate lives per-node (a scalar fold), but the answer must be reported for \emph{every} root -- the ``parent-side piece'' this chapter deliberately does not solve. Do the fold; hand off the every-root sweep to Ch.15.}
  \practice{Edge Weight Assignment}{Codeforces 1338B}
    {https://codeforces.com/problemset/problem/1338/B}{CF 1300 · L3}
    {Assign weights 1 or 2 to edges so that every leaf-to-leaf path has odd total weight; minimise the number of~$1$s. Colour by depth parity -- a scalar fold \coderef{c10-fold} with $S(v) = (\depth(v) \bmod 2)$ and $\Phi = $ XOR-into-parent. Everything else is counting.}
  \practice{Tree Shuffling}{Codeforces 1363E}
    {https://codeforces.com/problemset/problem/1363/E}{CF 1900 · L5}
    {Per-node swap costs; find the min total cost to sort a permutation living on a tree. Each node's state is a pair $(\text{cost if we shuffle in the subtree},\text{cost otherwise})$ -- the in/out idea of \thmref{c10-mis-tree} at a new combiner. $\bigO{n}$ once you see the reduction.}
  \practice{Collect Coins}{Codeforces 590E}
    {https://codeforces.com/problemset/problem/590/E}{CF 2100 · L5}
    {Long path on a grid with time-locked cells; classic DP with a two-row rolling state. Framed here as ``the same monotone-accumulator discipline as Shape B, but on a linear path instead of a tree.'' Deliberately not a tree-DP anchor -- spot the difference in the fold.}
\end{mixedset}

\chaptersection{The Summit -- CF 1076E Vasya and a Tree}

\begin{notebox}[Why this is the chapter top rope]
Sections 10.1--10.4 all fixed the root once and asked for a
\emph{read}. The Summit asks for a \emph{write pattern}: $m$ queries of
the form ``add $x$ to every node in $v$'s subtree at distance at most
$d$ from $v$'', then read the final value at every node. Naively each
query touches up to $n$ nodes -- $\bigO{nm}$ at $n,m\le 3\cdot 10^{5}$.
The fix is \thmref{c10-depth-xor}'s shape promoted from an XOR
accumulator to a per-depth \emph{lazy difference array}: on entry to a
node at depth $d$, add the deferred $\mathit{V}[d]$ into a running sum,
consume it, and apply every query at this node by writing $+x$ into the
running sum and $-x$ into $\mathit{V}[d+d_q+1]$ so the delta stops
applying past the correct depth. Siblings and cousins must see only the
queries relevant to their own subtree, so we undo $\mathit{V}[d]$ and
the running sum on exit. Every tool in the answer is already in
\coderef{c10-kit}: the DFS skeleton of \coderef{c10-fold}, the
monotone-accumulator insight of \coderef{c10-depth-hist}, and the
snapshot/undo discipline Ch.~12 introduced in \coderef{c09-iter}.
\end{notebox}

\solvedproblem{c10-summit}{Vasya and a Tree}{Codeforces 1076E}
  {https://codeforces.com/problemset/problem/1076/E}{CF 2000 · L5}
  {per-depth lazy difference on a fixed-root DFS, snapshot-and-undo accumulator, offline range updates on a subtree, long long throughout}

\begin{problemstatement}[Vasya and a Tree (samples verified against the live statement)]
A tree of $n$ vertices, vertex~$1$ is the root; the root has depth~$0$
and each child has depth one more than its parent. There are $m$
queries, each written $(v, d_q, x)$: add $x$ to every vertex in $v$'s
subtree at distance $\le d_q$ from $v$ (i.e., at absolute depth in
$[\depth(v),\,\depth(v) + d_q]$). Print the resulting value at every
vertex $1,2,\dots,n$ after all queries. Constraints:
$1\le n,m\le 3\cdot 10^{5}$; $0\le d_q\le 10^{9}$;
$1\le x\le 10^{9}$.\\
\textbf{Sample}: $n=5$, edges $1\!-\!2$, $1\!-\!3$, $2\!-\!4$,
$2\!-\!5$; $m=3$ queries $(1,1,1)$, $(2,0,10)$, $(4,10,100)$ $\Rightarrow$
output \texttt{1 11 1 100 0}. Query~1 adds $1$ to $\{1,2,3\}$ (nodes at
depth $\le 1$ in the whole tree); query~2 adds $10$ to $\{2\}$ (only
node~$2$, since $d_q=0$); query~3 adds $100$ to $\{4\}$ (node~$4$ is a
leaf). Total per node: $1$, $1+10$, $1$, $100$, $0$.
\end{problemstatement}

\begin{attemptgate}[Attempt it yourself first -- \lv{5}, give yourself 40 minutes]
Naive per-query BFS over $v$'s subtree is $\bigO{n\cdot m}$ and dies at
$3\cdot 10^{5}\cdot 3\cdot 10^{5}=9\cdot 10^{10}$. Two questions unlock
the linear shape: (i) what is the \emph{shape} of the update -- is it a
function of the absolute depth of the target, of the target's membership
in $v$'s subtree, or both? (ii) if it is both, which one does a DFS
traversing $v$'s subtree naturally track?
\peek{A DFS is already at $v$; every node in $v$'s subtree is visited
      in the window between $v$'s entry and $v$'s exit. Can you
      amortise the query onto that window?}
\nudge{Bucket queries by $v$. When you enter a node $u$ at depth $d$,
      every ancestor $v$ of $u$ that has a query $(v, d_q, x)$ with
      $d \le \depth(v) + d_q$ contributes $x$ to $\mathit{ans}[u]$.
      The active set changes by $+x$ on entering a query's $v$ and by
      $-x$ when the depth crosses the query's ceiling -- a lazy
      difference on the depth axis.}
\fullhint{Keep one running scalar $\mathit{cur}$ and one array
      $\mathit{V}[d]$ indexed by absolute depth. On entering $u$ at
      depth $d$: save $\mathit{cur}_0=\mathit{cur}$ and
      $\mathit{V}_0=\mathit{V}[d]$, then $\mathit{cur}\mathrel{+}=
      \mathit{V}[d]$, $\mathit{V}[d]=0$. For every query at $u$
      $(d_q, x)$: $\mathit{cur}\mathrel{+}=x$ and
      $\mathit{V}[d+d_q+1]\mathrel{-}=x$ (the delta expires past the
      budget). Read $\mathit{ans}[u]=\mathit{cur}$, recurse into
      children, and on exit restore $\mathit{cur}=\mathit{cur}_0$,
      $\mathit{V}[d]=\mathit{V}_0$. Total $\bigO{n+m}$ time and
      $\bigO{n+\max d_q\text{ clipped at }n}$ memory. Use \lstinline{long long}
      everywhere: $n\cdot \max x = 3\cdot 10^{14}$. Full build below.}
\end{attemptgate}

\solstep{Understanding the Problem}
Every query has two independent conditions on the target $u$: (a) $u$
is in $v$'s subtree (a \emph{DFS-window} condition -- $u$ is visited
between $v$'s entry and $v$'s exit) and (b) $\depth(u) \le
\depth(v) + d_q$ (a \emph{depth-ceiling} condition). A DFS already
gives you (a) for free, so the algorithm's job is to enforce (b) while
amortising the update. The lazy difference on the depth axis converts
``the update expires past a depth'' into two scalar writes on entry --
one to $+x$ the running sum, one to $-x$ the expiration bucket
$\mathit{V}[d+d_q+1]$.

\solstep{Key Observations}
\begin{observation}
Bucket queries by their $v$; the DFS visits $v$ exactly once, so the
$+x$ contribution can be applied once and inherited by every deeper
node through the running scalar $\mathit{cur}$.
\end{observation}
\begin{observation}
The depth ceiling is enforced by a single deferred write
$\mathit{V}[d+d_q+1]\mathrel{-}=x$: any node entering depth
$d+d_q+1$ (which includes $u$'s descendants one edge past the budget)
will consume the $-x$ and cancel the earlier $+x$. Nodes at
$\depth \le d_q + \depth(v)$ never enter the expiration bucket.
\end{observation}
\begin{observation}
Siblings and cousins must not see $u$'s contribution; the snapshot and
undo of $\mathit{cur}$ and $\mathit{V}[d]$ on $u$'s exit is what keeps
the subtree condition (a) intact. This is the \emph{only} place a
restore is legitimate in Ch.13's four-shape family.
\end{observation}
\begin{observation}
$\max d_q$ may be $10^{9}$, so $\mathit{V}[\cdot]$ is sized $\bigO{n}$
and index $d + d_q + 1$ is written \emph{only if it is $\le n$};
otherwise the delta never expires and there is nothing to cancel.
\end{observation}

\solstep{Think It Through With Me}
\begin{thinkbox}
\deadend{Per query, DFS $v$'s subtree and add $x$ to every target.}
{$\bigO{nm}$ worst case, $9\cdot 10^{10}$ at $n,m\le 3\cdot 10^{5}$ --
\tle. Amortisation must let all queries share one DFS.}
\deadend{Fenwick over tin/tout indices: subtree range add.}
{Works for condition (a) alone; the depth ceiling (b) makes the range
not contiguous in the Euler order -- you'd need a 2-D structure.
Ch.24's \chref{c18} machinery exists but this is a $4$-second problem
at $n=3\cdot 10^{5}$ and a 2-D BIT will \mle. Not the shape Ch.13
reaches for.}
\deadend{Euler-tour + per-node depth table, one pass per depth.}
{Correct but $\bigO{n + \height \cdot m}$; the depth axis is the same
size as the tree, so this collapses to the per-query DFS bound.}
\aha{One DFS, one running scalar, one $\mathit{V}[\cdot]$ per-depth
deferred-difference array, undo on exit. $\bigO{n+m}$ and \lstinline{long long}.}
\end{thinkbox}

\solstep{The Approach Ladder}
\begin{center}\renewcommand{\arraystretch}{1.3}
\begin{tabular}{@{}c p{0.235\linewidth} p{0.30\linewidth} l c@{}}
\toprule
\textbf{\#} & \textbf{Approach} & \textbf{Idea} & \textbf{Time} & \textbf{Verdict} \\
\midrule
1 & \lv{3} Per-query DFS & walk $v$'s subtree, apply $x$ up to depth & $O(nm)$ & \tle \\
2 & \lv{4} BIT on tin & subtree range add via Euler tour & $O((n+m)\log n)$ & \mle \\
3 & \lv{5} Per-depth lazy difference & one DFS, $\mathit{V}[d]$, snapshot-and-undo & $O(n+m)$ & \ac \\
\bottomrule
\end{tabular}\end{center}
Rung~1 is honest but pays per pair. Rung~2 uses Ch.18's Euler-tour
machinery but the ceiling condition fragments the update range, so a
1-D BIT does not suffice and the 2-D variant over-memoises. Rung~3 is
\emph{this} chapter's answer -- \thmref{c10-depth-xor}'s insight
promoted from XOR to a $+$-difference array, exactly the shape the
CF~570D listing sketched.

\solstep{Solution}
Read $n$ and the $n-1$ edges into an adjacency list. Read $m$ queries
and bucket them into $\mathit{qs}[v]$, each entry a pair
$(d_q, x)$. Allocate $\mathit{V}$ of size $n+2$ (clipping $d + d_q + 1$
at $n+1$ so the bucket never leaves the array) and $\mathit{ans}$ of
size $n+1$. Set $\mathit{cur}=0$. Run the DFS of \coderef{c10-summit}
from node~$1$ at depth~$0$: snapshot $\mathit{cur}$ and
$\mathit{V}[d]$, fold $\mathit{V}[d]$ into $\mathit{cur}$ and zero the
bucket, apply every query at $u$ ($\mathit{cur}\mathrel{+}=x$;
$\mathit{V}[d+d_q+1]\mathrel{-}=x$ when the index is in range), write
$\mathit{ans}[u]=\mathit{cur}$, recurse into children, then restore
$\mathit{cur}$ and $\mathit{V}[d]$. Print $\mathit{ans}[1],\dots,
\mathit{ans}[n]$. Every accumulator is \lstinline{long long}.

\solstep{Proof of Correctness}
\begin{proofbox}[\coderef{c10-summit} returns the aggregate of every query at every node]
\emph{Invariant.} Immediately before entering node $u$ at depth $d$,
$\mathit{cur}$ equals $\sum x_j$ over every query $j=(v_j,d_j,x_j)$
whose $v_j$ is an \emph{ancestor of $u$} (strict) with
$\depth(v_j) < d \le \depth(v_j) + d_j$; and $\mathit{V}[k]= -\sum_j x_j$
over queries $j$ whose ceiling $\depth(v_j)+d_j+1$ equals $k$ and whose
$v_j$ is an ancestor of $u$ (strict). No query with an off-path $v_j$
contributes, because the restore-on-exit of every node on the current
DFS path zeroed its writes.

\emph{Enter step.} On entering $u$: first fold $\mathit{V}[d]$ into
$\mathit{cur}$ and zero the bucket -- this cancels every $-x_j$ whose
ceiling is exactly $d$, correctly removing queries $j$ whose budget
ended at $d-1$; $\mathit{cur}$ is now $\sum x_j$ over every strict
ancestor $v_j$ with $d \le \depth(v_j) + d_j$. Then apply $u$'s own
queries: $\mathit{cur}\mathrel{+}=x$ and
$\mathit{V}[d+d_q+1]\mathrel{-}=x$ extends the invariant to include $u$
itself as an ancestor (of its subtree) with the correct ceiling.
$\mathit{ans}[u]$ is read after $u$'s own queries are folded in,
matching the definition (a query at $v$ covers $v$ itself since
$\depth(v) - \depth(v) = 0 \le d_q$).

\emph{Recurse.} Every descendant $w$ of $u$ at depth $d' \ge d$ enters
with $\mathit{cur}$ inherited and $\mathit{V}[\cdot]$ intact; the
invariant extends to $w$ by the same argument, restricted to $w$'s
ancestor set which is $u$'s ancestors plus $u$ itself.

\emph{Exit / undo.} Restoring $\mathit{cur}$ and $\mathit{V}[d]$
returns both to their pre-$u$ values, so siblings of $u$ see exactly
the strict-ancestor contributions they should and nothing from $u$'s
subtree.

\emph{Termination.} After the DFS, every $u$ has been entered and
exited exactly once, and $\mathit{ans}[u]$ was read at the correct
moment. $\blacksquare$ \qed
\end{proofbox}

\solstep{Implementation}
\begin{lstlisting}[caption={Vasya and a Tree -- per-depth lazy difference on a fixed-root DFS, snapshot-and-undo; long long throughout (Codeforces 1076E; O(n + m))}, label={lst:c10-summit}]
#include <bits/stdc++.h>
using namespace std;

int n, m;
vector<vector<int>> adj;
vector<vector<pair<int,long long>>> qs;   // qs[u] = list of (d_q, x)
vector<long long> V, ans;                 // V[d] = deferred delta for depth d
long long cur = 0;

void dfs(int u, int p, int d) {
    long long cur0 = cur;
    long long V0   = V[d];
    cur += V0;                            // consume the deferred delta
    V[d] = 0;                             // (bucket is now used for this subtree)
    for (auto& [dq, x] : qs[u]) {
        cur += x;
        int exp_d = d + dq + 1;           // absolute depth where the +x stops applying
        if (exp_d <= n) V[exp_d] -= x;    // clip at n+1 to keep V in range
    }
    ans[u] = cur;
    for (int v : adj[u]) if (v != p) dfs(v, u, d + 1);
    cur = cur0;                           // undo -- siblings must not see u's subtree contribution
    V[d]  = V0;
}

int main(){
    if (scanf("%d", &n) != 1) return 0;
    adj.assign(n + 1, {});
    for (int i = 0; i < n - 1; ++i){
        int u, v; scanf("%d %d", &u, &v);
        adj[u].push_back(v); adj[v].push_back(u);
    }
    qs.assign(n + 1, {});
    scanf("%d", &m);
    for (int i = 0; i < m; ++i){
        int v, dq; long long x; scanf("%d %d %lld", &v, &dq, &x);
        qs[v].push_back({dq, x});
    }
    V.assign(n + 2, 0); ans.assign(n + 1, 0);
    dfs(1, 0, 0);                          // root at 1, depth 0 (CF convention)
    for (int i = 1; i <= n; ++i) printf("%lld%c", ans[i], " \n"[i==n]);
    return 0;
}
\end{lstlisting}
\begin{notebox}[The four lines that are the whole idea]
\code{cur += V[d]; V[d] = 0;} consumes a deferred cancellation.
\code{cur += x; V[d + dq + 1] -= x;} applies a query and schedules its
expiry. \code{cur = cur0; V[d] = V0;} keeps siblings clean. Everything
else is reading input and printing the result. This is
\thmref{c10-depth-xor}'s ``monotone accumulator plus per-query
snapshot'' promoted from $\Phi=\oplus$ to $\Phi=+$.
\end{notebox}

\solstep{Trace}
\begin{tracebox}[Official Sample -- $n{=}5$, edges $1\!-\!2, 1\!-\!3, 2\!-\!4, 2\!-\!5$, three queries]
Root at~$1$ at depth~$0$. Depths: $1\!\to\!0$, $2\!\to\!1$, $3\!\to\!1$,
$4\!\to\!2$, $5\!\to\!2$. Children: $1\!\to\!\{2,3\}$, $2\!\to\!\{4,5\}$.
Bucketed queries: $\mathit{qs}[1]=\{(1,1)\}$, $\mathit{qs}[2]=\{(0,10)\}$,
$\mathit{qs}[4]=\{(10,100)\}$. Initially $\mathit{cur}=0$, $\mathit{V}[\cdot]=0$.
\\[3pt]
\textbf{Enter 1 at $d{=}0$}: $\mathit{cur}_0=0$, $\mathit{V}_0=0$.
$\mathit{cur}\mathrel{+}=\mathit{V}[0]=0$; $\mathit{V}[0]=0$. Apply query at~$1$:
$\mathit{cur}\mathrel{+}=1\Rightarrow \mathit{cur}=1$; $\mathit{V}[0+1+1]=\mathit{V}[2]\mathrel{-}=1\Rightarrow \mathit{V}[2]=-1$.
Read $\mathit{ans}[1]=\mathbf{1}$ \checkmark.\\
\textbf{Enter 2 at $d{=}1$}: $\mathit{cur}_0=1$, $\mathit{V}_1=\mathit{V}[1]=0$.
$\mathit{cur}\mathrel{+}=0\Rightarrow \mathit{cur}=1$; $\mathit{V}[1]=0$. Apply query at~$2$:
$\mathit{cur}\mathrel{+}=10\Rightarrow \mathit{cur}=11$; $\mathit{V}[1+0+1]=\mathit{V}[2]\mathrel{-}=10\Rightarrow \mathit{V}[2]=-11$.
Read $\mathit{ans}[2]=\mathbf{11}$ \checkmark.\\
\textbf{Enter 4 at $d{=}2$}: $\mathit{cur}_0=11$, $\mathit{V}_2=\mathit{V}[2]=-11$.
$\mathit{cur}\mathrel{+}=-11\Rightarrow \mathit{cur}=0$; $\mathit{V}[2]=0$. Apply query at~$4$:
$\mathit{cur}\mathrel{+}=100\Rightarrow \mathit{cur}=100$; $\mathit{V}[2+10+1]=\mathit{V}[13]$ -- index~$13$ exceeds $n{=}5$, clipped (no effect).
Read $\mathit{ans}[4]=\mathbf{100}$ \checkmark.\\
\textbf{Exit 4}: restore $\mathit{cur}=11$, $\mathit{V}[2]=-11$.\\
\textbf{Enter 5 at $d{=}2$}: $\mathit{cur}_0=11$, $\mathit{V}_2=\mathit{V}[2]=-11$.
$\mathit{cur}\mathrel{+}=-11\Rightarrow \mathit{cur}=0$; $\mathit{V}[2]=0$. No queries at~$5$.
Read $\mathit{ans}[5]=\mathbf{0}$ \checkmark.\\
\textbf{Exit 5}: restore $\mathit{cur}=11$, $\mathit{V}[2]=-11$.\\
\textbf{Exit 2}: restore $\mathit{cur}=1$, $\mathit{V}[1]=0$.\\
\textbf{Enter 3 at $d{=}1$}: $\mathit{cur}_0=1$, $\mathit{V}_1=0$.
$\mathit{cur}\mathrel{+}=0\Rightarrow \mathit{cur}=1$. No queries at~$3$.
Read $\mathit{ans}[3]=\mathbf{1}$ \checkmark.\\
\textbf{Exit 3}, \textbf{Exit 1}.\\[3pt]
\textbf{Printed}: \texttt{1 11 1 100 0} \checkmark --- matches the
official sample.\\
\textbf{Cross-check with a direct per-node enumeration}: node~$1$ gets
$+1$ (query~1); node~$2$ gets $+1$ (query~1) and $+10$ (query~2);
node~$3$ gets $+1$ (query~1); node~$4$ gets $+100$ (query~3); node~$5$
gets no query (query~1's ceiling $\depth(1)+1=1 < \depth(5)=2$; query~2
covers only $d_q=0$ from $2$). Matches \checkmark.
\end{tracebox}

\solstep{Complexity}
\begin{complexitybox}
Time $\bigO{n + m}$: every node contributes one enter and one exit
(constant work each), every query contributes two $\bigO{1}$ updates on
enter. Memory $\bigO{n + m}$: adjacency $\bigO{n}$, bucketed queries
$\bigO{m}$, $\mathit{V}$ array $\bigO{n}$, $\mathit{ans}$ array
$\bigO{n}$. At $n, m \le 3\cdot 10^{5}$: two $\mathit{cur}$/bucket
arrays and a $\bigO{m}$ query list is roughly 15\,MB -- comfortable
under CF's 256\,MB. On a legal chain of $3\cdot 10^{5}$, the recursive
DFS is $3\cdot 10^{5}$ frames deep -- swap for \coderef{c09-iter}'s
two-phase stack if your contest environment caps recursion at $\sim 10^{5}$.
Verdict: \ac.
\end{complexitybox}

\solstep{Verification and Edge Cases}
\begin{center}\renewcommand{\arraystretch}{1.2}
\begin{tabular}{@{}p{0.30\linewidth} p{0.62\linewidth}@{}}
\toprule
\textbf{Case} & \textbf{What the code does} \\
\midrule
$d_q$ exceeds $n$ & $\mathit{V}[d+d_q+1]$ clipped to $\le n+1$, so no expiration is scheduled; the $+x$ propagates to the deepest node in $v$'s subtree and no further, matching the ceiling. \\
$n = 1$ (single node) & One DFS frame; every query at~$1$ writes into $\mathit{cur}$, $\mathit{ans}[1]=\sum x$; the sample reduces to a scalar sum. \\
Two queries at the same $v$ & Both $+x$ applied, both expirations scheduled at possibly different depths -- the $\mathit{V}[\cdot]$ bucket is additive, no interference. \\
$x = 10^{9}$, $n=3\cdot 10^{5}$ & $\max \mathit{ans}[v]=3\cdot 10^{14}$, past \lstinline{int} by five orders; \lstinline{long long} throughout. \\
Zero-depth query ($d_q=0$) & Expiration scheduled at $d+1$ (the next depth), so only $v$ itself sees the $+x$ -- matches the sample's query~2 giving $+10$ only to node~$2$. \\
\bottomrule
\end{tabular}\end{center}

\solstep{Pitfalls and Variants}
\begin{warningbox}
Ways to break it: (i) \emph{no undo} -- the running sum leaks into
siblings and cousins; this is Shape~B's discipline inverted, and the
distinction is which of $\mathit{cur}$ (per-node scalar) versus
$\mathit{X}[h]$ (per-depth aggregate) you are maintaining. (ii)
\emph{int overflow} -- $\max \mathit{ans}[v]$ reaches $3\cdot 10^{14}$;
\lstinline{long long} is non-negotiable. (iii) \emph{indexing out of
bounds} -- $\mathit{V}$ must be sized $n+2$ and the expiration index
$d+d_q+1$ clipped at $n+1$. (iv) \emph{reading $\mathit{ans}[u]$ before
applying $u$'s own queries} -- the query at $v$ covers $v$ itself (the
sample's query~2 confirms $+10$ reaches node~$2$). \emph{Variants:}
replace $\mathit{V}[d+d_q+1] \mathrel{-}= x$ with $\mathit{V}[d + d_q + 1]
\mathrel{-}= x$ per child branch and you get a per-branch ceiling; the
``count queries whose target includes this node'' variant replaces $x$
by~$1$; and the general ``range-by-depth offline update'' family is a
recurring Ch.15 (\chref{c20}) pattern once the ceiling must hold for
every root.
\end{warningbox}

\begin{insightbox}[Key takeaways -- three things to remember]
\begin{enumerate}[nosep, leftmargin=1.6em]
  \item \emph{Amortise a query onto the DFS window.} Bucket by $v$;
        the subtree condition (a) is free; the depth condition (b)
        costs one $\mathit{V}[\cdot]$ slot per ceiling
        (\thmref{c10-depth-xor}, \probref{c10-summit}).
  \item \emph{Snapshot-and-undo is legitimate exactly when you are
        maintaining a per-node scalar.} Per-depth monotone accumulators
        (Shape~B, CF~570D) \emph{never} restore; per-node running sums
        (Summit, CF~1076E) \emph{must} restore -- the difference is
        what makes the pair ``same family, opposite disciplines'' click
        (\probref{c10-distk-solved} vs \probref{c10-summit}).
  \item \emph{Name the combiner before choosing the shape.} XOR, $+$,
        $\max$, $\sum$ --- the four chapter sections each instantiate
        one combiner on the same DFS; recognising the combiner is the
        whole decision (opener taxonomy chart, \coderef{c10-kit}).
\end{enumerate}
\end{insightbox}

\begin{hardeningbox}
\begin{itemize}[nosep, leftmargin=1.4em]
  \item \textbf{Generalises}: any offline ``range update on subtree
        plus a per-node or per-depth ceiling'' collapses to a
        DFS-embedded lazy-difference -- Ch.15's rerooting replays the
        same trick across all roots, and Ch.24's centroid decomposition
        runs it across path-length buckets (\chref{c20}, \chref{c18}).
  \item \textbf{The four listings of this chapter are one skeleton}:
        \coderef{c10-fold}, \coderef{c10-inout}, \coderef{c10-peak},
        and \coderef{c10-depth-hist} all share the shape ``one DFS, one
        per-frame state, one combiner''; the Summit is a fifth
        instantiation with a per-node scalar instead of a per-depth
        table (see the opener's taxonomy chart).
  \item \textbf{Alternatives}: an Euler-tour + 2-D BIT (Ch.18's
        \chref{c16}, Ch.24's \chref{c18}) solves this in
        $\bigO{(n+m)\log^2 n}$ but is over-memoised for the offline
        setting; a per-child subtree enumeration is $\bigO{nm}$ and
        TLEs at $n=m=3\cdot 10^{5}$ (\coderef{c10-kit}).
\end{itemize}
\end{hardeningbox}

\chaptersection{Tool Kit Inventory}

The machine assembled so far -- cumulative through Chapter~13.
Chapters~1--12's rows are carried forward; Chapter~13's rows are
appended. Pointers only, never prose.

\begin{center}
\small
\renewcommand{\arraystretch}{1.25}
\begin{tabular}{@{}p{0.22\linewidth} l c p{0.26\linewidth} p{0.18\linewidth}@{}}
\toprule
\textbf{Tool} & \textbf{Since} & \textbf{Level} & \textbf{Trigger phrase} & \textbf{Where} \\
\midrule
\multicolumn{5}{l}{\itshape Chapters 1--12 -- carried forward} \\
Binary-tree height/size recursion & 1 & \lv1 & ``compute at the node from its two children'' & \chref{c01} \\
Bottom-up DP pattern & 2 & \lv2 & ``a node folds its children's answers'' & \chref{c02} \\
Parent-guarded rooting / representation & 4 & \lv2 & ``root a general tree with a parent parameter'' & \chref{c04} \\
Measures DFS ($\depth/\height/\size$) & 12 & \lv2 & ``how deep / how tall / how many below'' & \coderef{c09-dfs3} \\
Iterative two-phase DFS & 12 & \lv3 & ``the chain may be $5\cdot10^{5}$ deep'' & \coderef{c09-iter} \\
DP diameter (join two arms at peak) & 12 & \lv3 & ``diameter that must carry weights / a predicate'' & \thmref{c09-dpjoin} \\
\midrule
\multicolumn{5}{l}{\itshape Chapter 13 -- new this chapter} \\
Subtree tuple fold & 10.1 & \lv3 & ``an answer per node from a fixed-size tuple'' & \coderef{c10-fold} \\
In/out pair (MIS / vertex cover) & 10.2 & \lv4 & ``take or leave $v$; forced out on children'' & \coderef{c10-inout}, \thmref{c10-mis-tree} \\
Peak-harvest identity & 10.3 & \lv4 & ``sum over all pairs with an additive $F$'' & \propref{c10-pairtop} \\
Distance-$k$ histogram merge & 10.3 & \lv4 & ``count pairs at distance $k$ in one DFS'' & \coderef{c10-peak} \\
Depth-indexed state (Shape A) & 10.4 & \lv4 & ``per-node histogram up to $k$'' & \coderef{c10-distk-hist}, \defref{c10-dist-hist} \\
Backtracking + monotone accumulator (Shape B) & 10.4 & \lv4 & ``offline per-depth aggregate with invertible $\Phi$'' & \coderef{c10-depth-hist}, \thmref{c10-depth-xor} \\
Per-depth lazy difference (Summit) & Summit & \lv5 & ``subtree + depth-ceiling range updates, offline'' & \coderef{c10-summit} \\
\bottomrule
\end{tabular}
\end{center}

\chaptersection{Related Topics}

\begin{itemize}[nosep, leftmargin=1.4em]
  \item \chref{c01}, \chref{c02} -- backward: the tuple fold of
        10.1 is Ch.~1's height recursion (\thmref{c01-height-recursion})
        promoted to a multi-coordinate state; Ch.2's bottom-up pattern
        (\chref{c02}) is the same shape with a specific tuple. Cited,
        never re-taught (\defref{c10-subtree-state}).
  \item \chref{c04} -- backward: rooting by an explicit parent
        parameter and the ``rooting + count-nodes-with-the-highest
        score'' family (\secref{c10-subtrees}'s LC~2049 back-cite) are
        Chapter~4's; this chapter takes the root as fixed and adds the
        combiner.
  \item \chref{c09} -- backward: Ch.12's \thmref{c09-dpjoin} is the
        $\Phi=\mathit{len}$-sum special case of \propref{c10-pairtop};
        Ch.12's $\depth/\height/\size$ table (\defref{c09-measures}) is
        the running example $T$ on which every Ch.13 hand trace sits.
  \item \chref{c27} -- forward: knapsack-over-children budget analysis
        and the formal $\sum_v \deg(v)\cdot k$ bound for the depth
        histogram are Chapter~14's; this chapter previews the shape
        without the proof (\secref{c10-distance}).
  \item \chref{c20} -- forward: rerooting \emph{as a theory} (any
        edge-transfer rule; answer for every root) is Chapter~15's;
        the ``parent-side piece'' this chapter deliberately does not
        maintain lives there (\probref{c10-summit}'s variants note
        names the exact extension).
  \item \chref{c11} -- forward: per-pair \emph{query} distance and the
        $\bigO{\log n}$ LCA machinery are Chapter~16's; this chapter
        discharges the forward references \propref{c10-pairtop},
        \secref{c10-distance}, and \exref{c10-k2} that Ch.~16's opener
        and chapter-end lean on.
  \item \chref{c28} -- forward: DSU-on-tree / small-to-large upgrades
        the histogram merge to $\bigO{n\log n}$ when $k$ degrades to
        $\Theta(\height)$; Ch.~13 previews this but the amortised
        analysis is Chapter~17's (\secref{c10-distance}).
  \item \chref{c18} -- forward: centroid decomposition recurses
        \emph{through} the peak-harvest identity of \propref{c10-pairtop}
        and gives an $\bigO{n\log n}$ alternative to the $\bigO{nk}$
        histogram; Chapter~24's (\secref{c10-harvest}'s forward cite).
\end{itemize}

\begin{lastword}[Last Word -- Chapter 13]
\emph{``Chapter~12 measured a tree; Chapter~13 makes it compute. The
root is fixed, the state is chosen, and one DFS is enough.''}

\medskip
The whole chapter is one discipline repeated four times: \emph{name a
state, name a combiner, run one DFS}. A tuple per node returns LC~333's
BST subtree check (\thmref{c10-subtree-template},
\probref{c10-subtree-solved}); a pair per node returns weighted MIS on
trees (\thmref{c10-mis-tree}, \probref{c10-inout-solved}); an arm
per peak returns the distance-$k$ census (\propref{c10-pairtop},
\probref{c10-harvest-solved}); a histogram per depth returns CF~570D's
palindrome queries in $\bigO{n+m}$ (\thmref{c10-depth-xor},
\probref{c10-distk-solved}). Ch.~12's \thmref{c09-dpjoin} was the
$\Phi=\max$-of-arm-sums instance; here it is one row in a table of
combiners. The Summit is the same DFS promoted from read to write --
per-depth lazy differences with a snapshot-and-undo discipline that is
\emph{opposite} to Shape~B's never-restore rule, because the object is
a per-node scalar rather than a per-depth aggregate.

\medskip
Carry three reflexes forward:
\begin{enumerate}[nosep, leftmargin=1.6em]
  \item \emph{Never pay per pair for a sum over pairs.} Harvest at the
        peak (\propref{c10-pairtop}), merge histograms in $\bigO{nk}$
        (\coderef{c10-peak}), or use an invertible accumulator in
        $\bigO{n+m}$ (\coderef{c10-depth-hist}). $\bigO{n^2}$ is
        always a shape mistake.
  \item \emph{The seed is the row that gets forgotten.}
        $\mathit{H}[u][0]=1$ is the arm of length~$0$; the ``peak
        equals one endpoint'' rows of \exref{c10-k2} ride on it, and
        Ch.~16's \exref{c10-k2} forward reference exists precisely
        because those rows are easy to skip.
  \item \emph{Name the combiner before choosing the shape.} Sum $\to$
        histogram merge or lazy difference; XOR (or any group op) $\to$
        Shape~B; $\max$ over additive arms $\to$ peak-harvest; a fixed
        boolean decision per node $\to$ the in/out pair. The deciding
        word is in the problem text, and this chapter's four sections
        are the four answers.
\end{enumerate}

\medskip
And three habits, drilled until boring: keep the Ch.~12 measures DFS
within arm's reach because every DP state is that DFS with a heavier
frame; check the length-$0$ row on every distance-indexed count before
you ship; and when a listing has to restore an accumulator, justify
whether it is per-node scalar (Summit -- restore) or per-depth aggregate
(CF~570D -- never restore) in one sentence (\coderef{c10-kit}).

\medskip
\textbf{Next:} \chref{c27} takes the histogram of \secref{c10-distance}
and asks the harder budget question -- \emph{when} does the
$\bigO{nk}$ merge amortise, and \emph{what} does a knapsack-over-
children state cost across a tree of non-constant degree? Ch.~13 hands
Ch.~14 four shapes and a running example; Ch.~14 gives you a
$\sum_v \deg(v) \cdot k$ theorem, a small-to-large preview, and a
family of profile DPs where the state grows past a fixed tuple. The
cookbook is open -- from here, the analysis becomes the point.
\end{lastword}
